\ifdefined\gkver\else\def\gkver{student}\fi
\documentclass[\gkver]{gaokaozhenti}
\begin{document}

\chapter{1997年全国}

\begin{enumerate}
\item
%% number: 1
%% paperName: 1997年全国高考第 1 题 3 分
%% typeId: 1
%% type: 单选题
%% chapter: 原子和原子核
%% point: 原子核式结构
%% method: 无
%% score: 3
%% degree: 850
%% duplicateId: 0
%% body:
在卢瑟福的 $\alpha$ 粒子散射实验中，有少数 $\alpha$ 粒子发生大角度偏转，其原因是\xzanswer{A}
\fourchoices[answer=A]
{原子的正电荷和绝大部分质量集中在一个很小的核上}
{正电荷在原子中是均匀分布的}
{原子中存在着带负电的电子}
{原子只能处于一系列不连续的能量状态中}

%% answer: A
%% memo:
\memoanswer{
由于原子的全部正电荷和绝大部分质量都集中在一个很小的核上，使得整个原子内部十分“空旷”，当 $\alpha$ 粒子穿过原子时，接近或正碰原子核的机会极少，所以只能有少数的 $\alpha$ 粒子发生大角度偏转，故 A 正确。
}

\item
%% number: 2
%% paperName: 1997年全国高考第 2 题 3 分
%% typeId: 1
%% type: 单选题
%% chapter: 运动和力
%% point: 动量定理
%% method: 无
%% score: 3
%% degree: 650
%% duplicateId: 0
%% body:
质量为 $m$ 的钢球自高处落下，以速率 $v_{1}$ 碰地，竖直向上弹回，碰撞时间极短，离地的速率为 $v_{2}$。在碰撞过程中，地面对钢球的冲量的方向和大小为\xzanswer{D}
\fourchoices[answer=D]
{向下，$m(v_{1}-v_{2})$}
{向下，$m(v_{1}+v_{2})$}
{向上，$m(v_{1}-v_{2})$}
{向上，$m(v_{1}+v_{2})$}

%% answer: D
%% memo:
\memoanswer{
物体以竖直速度 $v_{1}$ 与地面碰撞，而又以 $v_{2}$ 的速度反弹。设垂直地面向上的方向为正方向，对钢球应用动量定理得 $Nt-mgt=mv_{2}-(-mv_{1})=mv_{2}+mv_{1}$，由于碰撞时间极短，则 $mgt$ 可忽略不计，所以 $Nt=m(v_{2}+v_{1})$，即弹力的冲量方向向上，大小为 $m(v_{2}+v_{1})$，故 D 正确。
}

\item
%% number: 3
%% paperName: 1997年全国高考第 3 题 3 分
%% typeId: 1
%% type: 单选题
%% chapter: 运动和力
%% point: 牛顿第二定律
%% method: 无
%% score: 3
%% degree: 650
%% duplicateId: 0
%% body:
质量为 $M$ 的木块位于粗糙水平桌面上，若用大小为 $F$ 的水平恒力拉木块，其加速度为 $a$。当拉力方向不变，大小变为 $2F$ 时，木块的加速度为 $a'$，则\xzanswer{C}
\fourchoices[answer=C]
{$a'=a$}
{$a'<2a$}
{$a'>2a$}
{$a'=2a$}

%% answer: C
%% memo:
\memoanswer{
因木块在粗糙水平桌面上运动，受到摩擦力作用，由牛顿第二定律有 $F-F_{f}=ma$ ①，当拉力变成 $2F$ 时，有 $2F-F_{f}=ma'$ ②，由②式减①式得 $F=m(a'-a)$ ③，将①③式相比较，因 $F>F-F_{f}$，即 $m(a'-a)>ma$，得 $a'>2a$，故 C 正确。
}

\item
%% number: 4
%% paperName: 1997年全国高考第 4 题 3 分
%% typeId: 1
%% type: 单选题
%% chapter: 交流电
%% point: 变压器
%% method: 无
%% score: 3
%% degree: 650
%% duplicateId: 0
%% body:
A、B 两电路中，当 a、b 两端与 e、f 两端分别加上 $220\UV$ 的交流电压时，测得 c、d 间与 g、h 间的电压均为 $110\UV$。若分别在 c、d 两端与 g、h 两端加上 $110\UV$ 的交流电压，则 a、b 间与 e、f 间的电压分别为\xzanswer{B}
\onepicture[w=7cm]{figs/04.png}
\fourchoices[answer=B]
{$220\UV$，$220\UV$}
{$220\UV$，$110\UV$}
{$110\UV$，$110\UV$}
{$220\UV$，0}

%% answer: B
%% memo:
\memoanswer{
首先要搞清楚变压器和滑动变阻器在改变电压原理上的本质不同：对于变压器，a、b 与 c、d 间的电压比总是等于它们间线圈的匝数比，与哪一个是原线圈无关，故 a、b 间接 $220\UV$ 交变电压，c、d 间电压为 $110\UV$；c、d 间改接 $110\UV$ 交变电压，则 a、b 间应输出电压 $220\UV$；而对滑动变阻器，当 e、f 间接 $220\UV$ 电压时，电阻的 e、g 与 f、g 部分串联，g、h 间电压仅是 f、g 部分电阻的电压，当 g、h 间接 $110\UV$ 电压时，由于 e、g 部分无电流，e、g 两点等电势，故 e、f 间电压等于 g、h 间电压，故 B 正确。
}

\item
%% number: 5
%% paperName: 1997年全国高考第 5 题 3 分
%% typeId: 1
%% type: 单选题
%% chapter: 光学
%% point: 光的干涉
%% method: 无
%% score: 3
%% degree: 650
%% duplicateId: 0
%% body:
在双缝干涉实验中，以白光为光源，在屏幕上观察到了彩色干涉条纹，若在双缝中的一缝前放一红色滤光片（只能透过红光），另一缝前放一绿色滤光片（只能透过绿光），这时\xzanswer{C}
\fourchoices[answer=C]
{只有红色和绿色的双缝干涉条纹，其它颜色的双缝干涉条纹消失}
{红色和绿色的双缝干涉条纹消失，其它颜色的双缝干涉条纹依然存在}
{任何颜色的双缝干涉条纹都不存在，但屏上仍有光亮}
{屏上无任何光亮}

%% answer: C
%% memo:
\memoanswer{
两列光波发生干涉的条件之一是频率相等，利用双缝将一束光分成能够发生干涉的两束光，在光屏上形成干涉条纹，但分别用绿滤光片和红滤光片挡住两条缝后，红光和绿光频率不等，不能发生干涉，因此屏上不会出现干涉条纹，故 C 正确。
}

\item
%% number: 6
%% paperName: 1997年全国高考第 6 题 5 分
%% typeId: 2
%% type: 多选题
%% chapter: 原子和原子核
%% point: 人工核反应
%% method: 无
%% score: 5
%% degree: 650
%% duplicateId: 0
%% body:
在下列核反应方程中，x 代表质子的方程是\xzanswer{BC}
\fourchoices[answer=BC]
{\ce{^{27}_{13}Al + ^{4}_{2}He -> ^{30}_{15}P + x}}
{\ce{^{14}_{7}N + ^{4}_{2}He -> ^{17}_{8}O + x}}
{\ce{^{2}_{1}H + \gamma{} -> ^{1}_{0}n + x}}
{\ce{^{3}_{1}H + x -> ^{4}_{2}He + ^{1}_{0}n}}

%% answer: BC
%% memo:
\memoanswer{
解法一：将质子写成 \ce{^{1}_{1}H} 代替式中 x，然后由质量数守恒和电荷数守恒逐项检验，同时满足质量数和核电荷数守恒时，即为正确选项。

解法二：先由质量数守恒算出 x 的质量数，再由电荷数守恒算出 x 的电荷数，若两者均等于 1，x 便是质子，否则 x 代表其他粒子。如 A 中 x 的质量数为 $27+4-30=1$，电荷数为 $13+2-15=0$，所以 x 为中子（\ce{^{1}_{0}n}）；B 式中 x 的质量数为 $14+4-17=1$，电荷数为 $7+2-8=1$，所以 x 为质子；同理分析 C、D 选项可知，正确的答案为 BC。
}

\item
%% number: 7
%% paperName: 1997年全国高考第 7 题 5 分
%% typeId: 1
%% type: 单选题
%% chapter: 光学
%% point: 光的折射
%% method: 无
%% score: 5
%% degree: 650
%% duplicateId: 0
%% body:
光线在玻璃和空气的分界面上发生全反射的条件是\xzanswer{B}
\fourchoices[answer=B]
{光从玻璃射到分界面上，入射角足够小}
{光从玻璃射到分界面上，入射角足够大}
{光从空气射到分界面上，入射角足够小}
{光从空气射到分界面上，入射角足够大}

%% answer: B
%% memo:
\memoanswer{
发生全反射的条件为光从光密介质射入光疏介质，即光从玻璃射到空气中，且入射角须大于临界角，即要求入射角足够大，故 B 正确，A、C、D 错误。
}

\item
%% number: 8
%% paperName: 1997年全国高考第 8 题 5 分
%% typeId: 2
%% type: 多选题
%% chapter: 气体性质
%% point: 分子力
%% method: 无
%% score: 5
%% degree: 650
%% duplicateId: 0
%% body:
在下列叙述中，正确的是\xzanswer{AD}
\fourchoices[answer=AD]
{物体的温度越高，分子热运动越剧烈，分子平均动能越大}
{布朗运动就是液体分子的热运动}
{对一定质量的气体加热，其内能一定增加}
{分子间的距离 $r$ 存在某一值 $r_{0}$，当 $r<r_{0}$ 时，斥力大于引力，当 $r>r_{0}$ 时，斥力小于引力}

%% answer: AD
%% memo:
\memoanswer{
温度是分子平均动能的标志，温度升高，分子平均动能增大，分子热运动越剧烈，故 A 正确；布朗运动是指悬浮在液体中花粉（微小颗粒）的无规则运动，不是指液体分子的运动，故 B 错误；由热力学第一定律 $\Delta U=W+Q$ 知，对一定质量的气体加热（$Q$ 为正），但因做功情况 $W$ 未知，其内能的变化情况不能确定，故 C 错误；由分子力的特点与分子力随分子间距离变化的规律知 D 正确。
}

\item
%% number: 9
%% paperName: 1997年全国高考第 9 题 5 分
%% typeId: 2
%% type: 多选题
%% chapter: 力物体的平衡
%% point: 三力平衡
%% method: 无
%% score: 5
%% degree: 650
%% duplicateId: 0
%% body:
如图 \ref{1997全国09} 所示，图中重物的质量为 $m$，轻细线 AO 和 BO 的 A、B 端是固定的，平衡时 AO 是水平的，BO 与水平面的夹角为 $\theta$，AO 的拉力 $F_{1}$ 和 BO 的拉力 $F_{2}$ 的大小是\xzanswer{BD}
\onepicture[num=9, showlabel=false, label=1997全国09, w=4cm]{\includesvg{figs/09}}
\fourchoices[answer=BD]
{$F_{1}=mg\cos\theta$}
{$F_{1}=mg\cot\theta$}
{$F_{2}=mg\sin\theta$}
{$F_{2}=\frac{mg}{\sin\theta}$}

%% answer: BD
%% memo:
\memoanswer{
把 $F_{2}$ 分解为水平方向和竖直方向，由共点力的平衡得 $F_{2}\cos\theta=F_{1}$，$F_{2}\sin\theta=mg$，解得 $F_{2}=\frac{mg}{\sin\theta}$，$F_{1}=\frac{mg}{\tan\theta}$，故 B、D 正确。

\onepicture[w=4cm]{figs/09_an.jpg}
}

\item
%% number: 10
%% paperName: 1997年全国高考第 10 题 5 分
%% typeId: 1
%% type: 单选题
%% chapter: 交流电
%% point: LC振荡回路
%% method: 无
%% score: 5
%% degree: 650
%% duplicateId: 0
%% body:
为了增大 LC 振荡电路的固有频率，下列办法中可采取的是\xzanswer{D}
\fourchoices[answer=D]
{增大电容器两极板的正对面积并在线圈中放入铁芯}
{减小电容器两极板的距离并增加线圈的匝数}
{减小电容器两极板的距离并在线圈中放入铁芯}
{减小电容器两极板的正对面积并减少线圈的匝数}

%% answer: D
%% memo:
\memoanswer{
LC 振荡电路的固有频率为 $f=\frac{1}{2\pi\sqrt{LC}}$，若增大电容并增大电感，固有频率减小，故 A 错误；若增大电容并增大电感，固有频率减小，故 B 错误；若增大电容并增大电感，固有频率减小，故 C 错误；减小电容并减小电感，固有频率增大，故 D 正确。
}

\item
%% number: 11
%% paperName: 1997年全国高考第 11 题 5 分
%% typeId: 2
%% type: 多选题
%% chapter: 机械波
%% point: 波的图象
%% method: 无
%% score: 5
%% degree: 650
%% duplicateId: 0
%% body:
简谐横波某时刻的波形图线如图所示。由此图可知\xzanswer{BD}
\onepicture[w=7cm]{figs/11.png}
\fourchoices[answer=BD]
{若质点 a 向下运动，则波是从左向右传播的}
{若质点 b 向上运动，则波是从左向右传播的}
{若波从右向左传播，则质点 c 向下运动}
{若波从右向左传播，则质点 d 向上运动}

%% answer: BD
%% memo:
\memoanswer{
利用同侧法判断分析，若质点 a 向下运动，此波应向左传播，故 A 错误；若质点 b 向上运动，此波应向右传播，故 B 正确；若波向左传播，质点 c 应向上运动，故 C 错误；若波向左传播，质点 d 应向上运动，故 D 正确。
}

\item
%% number: 12
%% paperName: 1997年全国高考第 12 题 5 分
%% typeId: 2
%% type: 多选题
%% chapter: 电路
%% point: 电容
%% method: 无
%% score: 5
%% degree: 650
%% duplicateId: 0
%% body:
如图所示的电路中，电源的电动势恒定，要想使灯泡变暗，可以\xzanswer{AD}
\onepicture[w=6cm]{figs/12.png}
\fourchoices[answer=AD]
{增大 $R_{1}$}
{减小 $R_{1}$}
{增大 $R_{2}$}
{减小 $R_{2}$}

%% answer: AD
%% memo:
\memoanswer{
灯泡 R 与电阻 $R_{2}$ 并联，再与电阻 $R_{1}$ 串联分压，然后逐项分析：若增大 $R_{1}$，$R_{2}$ 上分压减小，灯泡上分压减小，所以灯泡变暗，故 A 正确；减小 $R_{1}$，$R_{1}$ 上分压减小，灯泡上分压增大，所以灯泡变亮，故 B 错误；增大 $R_{2}$，灯泡 R 与电阻 $R_{2}$ 并联电阻增大，灯泡上分压增大，所以灯泡变亮，故 C 错误；减小 $R_{2}$，灯泡 R 与电阻 $R_{2}$ 并联电阻减小，灯泡上分压减小，所以灯泡变暗，故 D 正确。
}

\item
%% number: 13
%% paperName: 1997年全国高考第 13 题 5 分
%% typeId: 2
%% type: 多选题
%% chapter: 电磁感应
%% point: 自感与互感，涡流
%% method: 无
%% score: 5
%% degree: 650
%% duplicateId: 0
%% body:
如图所示的电路中，A$_{1}$ 和 A$_{2}$ 是完全相同的灯泡，线圈 L 的电阻可以忽略。下列说法中正确的是\xzanswer{AD}
\onepicture[w=7cm]{figs/13.png}
\fourchoices[answer=AD]
{合上开关 K 接通电路时，A$_{2}$ 先亮，A$_{1}$ 后亮，最后一样亮}
{合上开关 K 接通电路时，A$_{1}$ 和 A$_{2}$ 始终一样亮}
{断开开关 K 切断电路时，A$_{2}$ 立刻熄灭，A$_{1}$ 过一会儿才熄灭}
{断开开关 K 切断电路时，A$_{1}$ 和 A$_{2}$ 都要过一会儿才熄灭}

%% answer: AD
%% memo:
\memoanswer{
自感线圈具有阻碍电流变化的作用，当电流增加时，它阻碍电流增加；当电流减小时，它阻碍电流减小，但阻碍并不是阻止。闭合开关，L 中电流从无到有，L 将阻碍这一变化，使 L 中电流不能迅速增大，而无电感的电路，电流能够瞬时达到稳定值，故 $A_{1}$ 灯后亮，$A_{2}$ 灯先亮，最后两灯电流相等，一样亮；断开开关时，L 中产生自感电流与自身原来的电流方向相同，自感电流通过 $A_{1}$、$A_{2}$，使 $A_{1}$、$A_{2}$ 过一会熄灭，故 A、D 正确。
}

\item
%% number: 14
%% paperName: 1997年全国高考第 14 题 5 分
%% typeId: 1
%% type: 单选题
%% chapter: 电路
%% point: 故障分析
%% method: 无
%% score: 5
%% degree: 650
%% duplicateId: 0
%% body:
在如图所示电路的三根导线中，有一根是断的，电源、电阻器 $R_{1}$、$R_{2}$ 及另外两根导线都是好的。为了查出断导线，某学生想先将多用表的红笔连接在电源的正极 a，再将黑表笔分别连接在电阻器 $R_{1}$ 的 b 端和 $R_{2}$ 的 c 端，并观察多用表指针的示数。在下列选档中，符合操作规程的是\xzanswer{A}
\onepicture[w=7cm]{figs/14.png}
\fourchoices[answer=A]
{直流 $10\UV$ 档}
{直流 $0.5\UA$ 档}
{直流 $2.5\UV$ 档}
{欧姆档}

%% answer: A
%% memo:
\memoanswer{
若操作时选用直流电压挡，由于电源电压是 $6\UV$，检查电路时只能用直流 $10\UV$ 挡，而不能用直流 $2.5\UV$ 挡，故 A 正确，C 错误；若操作时选直流电流挡，经估算通过电表的电流可能会达到 $1.2\UA$，它超过了 $0.5\UA$，这是不安全的，故 B 错误；用欧姆表去检查电路时，必须将电路中的电源断开，否则会损坏多用电表，况且电路中的电源与表头构成回路，无法判断出哪根导线是断的，故 D 错误。
}

\item
%% number: 15
%% paperName: 1997年全国高考第 15 题 5 分
%% typeId: 4
%% type: 实验题
%% chapter: 直线运动
%% point: 游标卡尺和螺旋测微仪
%% method: 无
%% score: 5
%% degree: 850
%% duplicateId: 0
%% body:
一游标卡尺的主尺最小分度为 $1\Umm$，游标上有 10 个小等分间隔，现用此卡尺来测量工件的直径，如图所示。该工件的直径为\tkanswer[2.5cm]{$29.80\Umm$}。
\onepicture[align=right, w=8cm]{figs/15.png}

%% answer: $29.80\Umm$
\jdanswer{
$29.80\Umm$
}

%% memo:
\memoanswer{
主尺读数是 $29\Umm$，游标上每一分度是 $0.1\Umm$，游标上第 8 条刻度线与主尺上毫米刻度线对齐，所以读数是 $29\Umm+8\times0.1\Umm=29.8\Umm$。
}

\item
%% number: 16
%% paperName: 1997年全国高考第 16 题 6 分
%% typeId: 4
%% type: 实验题
%% chapter: 气体性质
%% point: 验证玻意耳定律
%% method: 无
%% score: 6
%% degree: 650
%% duplicateId: 0
%% body:
下列给出的器材中，哪些是“验证玻-马定律实验”所必需的，把这些器材前面的字母填在横线上。

（A）带有刻度的注射器　　（B）刻度尺

（C）弹簧秤　　（D）钩码若干个

答：\tkanswer[2cm]{ABCD}。

实验读数过程中，不能用手握住注射器，这是为了\tkanswer[3cm]{保持气体的温度恒定}。

用橡皮帽封住注射器小孔，这是为了\tkanswer[3cm]{保持气体的质量不变}。

%% answer: ABCD；保持气体的温度恒定；保持气体的质量不变
\jdanswer{
\begin{enumerate}
	\item
	ABCD
	\item
	保持气体的温度恒定
	\item
	保持气体的质量不变
\end{enumerate}
}

%% memo:
\memoanswer{
ABCD；保持气体的温度恒定；保持气体的质量不变。
}

\item
%% number: 17
%% paperName: 1997年全国高考第 17 题 6 分
%% typeId: 4
%% type: 实验题
%% chapter: 电路
%% point: 测电动势与内阻
%% method: 无
%% score: 6
%% degree: 650
%% duplicateId: 0
%% body:
某电压表的内阻在 $20\sim50\UkO$ 之间，现要测量其内阻，实验室提供下列可选用的器材：

待测电压表 V（量程 $3\UV$）

电流表 $A_{1}$（量程 $200\UuA$）

电流表 $A_{2}$（量程 $5\UmA$）

电流表 $A_{3}$（量程 $0.6\UA$）

滑动变阻器 $R$（最大阻值 $1\UkO$）

电源 $\varepsilon$（电动势 $4\UV$）

电键 K。
\begin{enumerate}
	\item
	所提供的电流表中，应选用\tkanswer[2cm]{$A_{1}$}（填写字母代号）。
	\item
	为了尽量减小误差，要求测多组数据。试在方框中画出符合要求的实验电路图（其中电源和电键及其连线已画出）。
\end{enumerate}
\onepicture[align=right, w=6cm]{figs/17.png}

%% answer: （1）$A_{1}$；（2）如图
\jdanswer{
\begin{enumerate}
	\item
	$A_{1}$
	\item
	如图
\end{enumerate}
}

%% memo:
\memoanswer{
由于电源电动势 $E=4\UV$，电压表内阻最大 $50\UkO$，故电流最小值 $I=\frac{E}{R_{V}}=\frac{4}{50\times10^{3}}\UA=80\UuA$，$A_{2}$、$A_{3}$ 量程都太大，故选 $A_{1}$。由于 $R_{V}$ 在 $20\sim50\UkO$ 之间，远大于滑动变阻器的阻值 $1\UkO$，故应连成分压电路，电路图如图所示。

\onepicture[align=right, w=5cm]{figs/17_an.jpg}
}

\item
%% number: 18
%% paperName: 1997年全国高考第 18 题 5 分
%% typeId: 3
%% type: 填空题
%% chapter: 磁场
%% point: 带电粒子在磁场中的运动
%% method: 无
%% score: 5
%% degree: 650
%% duplicateId: 0
%% body:
如图 \ref{1997全国18}，在 $x$ 轴的上方（$y\geq0$）存在着垂直于纸面向外的匀强磁场，磁感应强度为 $B$。在原点 O 有一个离子源向 $x$ 轴上方的各个方向发射出质量为 $m$、电量为 $q$ 的正离子，速率都为 $v$。对那些在 $xy$ 平面内运动的离子，在磁场中可能到达的最大 $x=$\tkanswer[2.5cm]{$\frac{2mv}{qB}$}，最大 $y=$\tkanswer[2.5cm]{$\frac{2mv}{qB}$}。
\onepicture[num=18, showlabel=false, label=1997全国18, align=right, w=5cm]{\includesvg{figs/18}}

%% answer: $\frac{2mv}{qB}$，$\frac{2mv}{qB}$
\jdanswer{
$\frac{2mv}{qB}$，$\frac{2mv}{qB}$
}

%% memo:
\memoanswer{
由左手定则判断，从 O 点射出沿 $y$ 轴正方向运动的离子受洛伦兹力的方向沿 $Ox$ 正方向，正离子到达 $x$ 轴上的 $x$ 坐标最大，最大值为轨迹半径的 2 倍，由 $R=\frac{mv}{qB}$ 得 $x=2R=\frac{2mv}{qB}$；同理，从 O 点沿 $x$ 轴负方向射出的正离子到达 $y$ 轴的 $y$ 坐标最大，最大值也为 $2R$，即 $y=\frac{2mv}{qB}$。
}

\item
%% number: 19
%% paperName: 1997年全国高考第 19 题 5 分
%% typeId: 3
%% type: 填空题
%% chapter: 电场
%% point: 带电粒子的圆周运动
%% method: 无
%% score: 5
%% degree: 450
%% duplicateId: 0
%% body:
质量为 $m$、电量为 $q$ 的质点，在静电力作用下以恒定速率 $v$ 沿圆弧从 A 点运动到 B 点，其速度方向改变的角度为 $\theta$（弧度），AB 弧长为 $s$。则 A，B 两点间的电势差 $\varphi_{A}-\varphi_{B}=$\tkanswer[1.5cm]{0}，AB 弧中点场强大小 $E=$\tkanswer[3cm]{$\frac{mv^{2}\theta}{qs}$}。

%% answer: 0，$\frac{mv^{2}\theta}{qs}$
\jdanswer{
$0$，$\frac{mv^{2}\theta}{qs}$
}

%% memo:
\memoanswer{
质点以恒定的速率沿圆弧从 A 点运动到 B 点，电场力做功为零，则 A、B 两点间的电势差 $\varphi_{A}-\varphi_{B}=0$。设场源电荷的电量为 $Q$，质点做圆周运动的半径为 $r$，则 AB 的弧长 $s=\theta r$，静电力是质点做圆周运动的向心力，有 $k\frac{Qq}{r^{2}}=m\frac{v^{2}}{r}$，AB 弧中点的场强大小为 $E=\frac{kQ}{r^{2}}$，联立解得 $E=\frac{mv^{2}\theta}{qs}$。
}

\item
%% number: 20
%% paperName: 1997年全国高考第 20 题 5 分
%% typeId: 3
%% type: 填空题
%% chapter: 曲线运动
%% point: 人造地球卫星
%% method: 近似与估算
%% score: 5
%% degree: 450
%% duplicateId: 0
%% body:
已知地球半径约为 $6.4\times10^{6}\Um$，又知月球绕地球的运动可近似看作匀速圆周运动，则可估算出月球到地心的距离约为\tkanswer[2.8cm]{$4\times10^{8}\Um$}。（结果只保留一位有效数字）

%% answer: $4\times10^{8}\Um$
\jdanswer{
$4\times10^{8}\Um$
}

%% memo:
\memoanswer{
地球对月球的万有引力提供向心力，由牛顿第二定律得 $G\frac{Mm}{r^{2}}=m\frac{4\pi^{2}}{T^{2}}r$，又在地球表面 $G\frac{Mm}{R^{2}}=mg$，联立解得 $r=\sqrt[3]{\frac{gR^{2}T^{2}}{4\pi^{2}}}$，代入数据得 $r\approx4\times10^{8}\Um$。
}

\item
%% number: 21
%% paperName: 1997年全国高考第 21 题 5 分
%% typeId: 3
%% type: 填空题
%% chapter: 曲线运动
%% point: 竖直平面圆周运动
%% method: 无
%% score: 5
%% degree: 450
%% duplicateId: 0
%% body:
一内壁光滑的环形细圆管，位于竖直平面内，环的半径为 $R$（比细管的半径大得多）。在圆管中有两个直径与细管内径相同的小球（可视为质点）。A 球的质量为 $m_{1}$，B 球的质量为 $m_{2}$。它们沿环形圆管顺时针运动，经过最低点时的速度都为 $v_{0}$。设 A 球运动到最低点时，B 球恰好运动到最高点，若要此时两球作用于圆管的合力为零，那么 $m_{1}$、$m_{2}$、$R$ 与 $v_{0}$ 应满足的关系式\tkanswer[5cm]{$\frac{(m_{1}-m_{2})v_{0}^{2}}{R}+(m_{1}+5m_{2})g=0$}。

%% answer: $\frac{(m_{1}-m_{2})v_{0}^{2}}{R}+(m_{1}+5m_{2})g=0$
\jdanswer{
$\frac{(m_{1}-m_{2})v_{0}^{2}}{R}+(m_{1}+5m_{2})g=0$
}

%% memo:
\memoanswer{
A 球做圆周运动到最低点时，受重力 $m_{1}g$，圆管对其施加的向上的弹力 $N_{1}$，由牛顿第二定律得 $N_{1}-m_{1}g=\frac{m_{1}v_{0}^{2}}{R}$，解得 $N_{1}=m_{1}g+\frac{m_{1}v_{0}^{2}}{R}$；由牛顿第三定律知，A 球作用于管的弹力 $N_{1}'$ 方向向下，大小为 $N_{1}'=N_{1}=m_{1}g+\frac{m_{1}v_{0}^{2}}{R}$。若要此时两球作用于圆管的合力为零，则 B 球对管应施加向上的作用力，而 B 球则受向下的弹力 $N_{2}$ 和重力 $m_{2}g$，对 B 球，由牛顿第二定律得 $N_{2}+m_{2}g=\frac{m_{2}v^{2}}{R}$，解得 $N_{2}=\frac{m_{2}v^{2}}{R}-m_{2}g$。B 球对管的作用力 $N_{2}'=N_{2}$，两球作用于管的合力为 0，有 $N_{1}'=N_{2}'$，即 $m_{1}g+\frac{m_{1}v_{0}^{2}}{R}=\frac{m_{2}v^{2}}{R}-m_{2}g$ ①。B 球由最低点向最高点运动过程中，只有重力对其做功，机械能守恒，对 B 球，由机械能守恒得 $\frac{1}{2}m_{2}v_{0}^{2}=m_{2}g\cdot2R+\frac{1}{2}m_{2}v^{2}$ ②。联立①②两式解得 $(m_{1}-m_{2})\frac{v_{0}^{2}}{R}+(m_{1}+5m_{2})g=0$。
}

\item
%% number: 22
%% paperName: 1997年全国高考第 22 题 9 分
%% typeId: 6
%% type: 计算题
%% chapter: 光学
%% point: 透镜
%% method: 无
%% score: 9
%% degree: 750
%% duplicateId: 0
%% body:
有一个焦距为 $36\Ucm$ 的凸透镜，在主轴上垂直放置一支蜡烛，得到一个放大率为 4 的虚像。如果想得到放大率为 4 的实像，蜡烛应向哪个方向移动？移动多少？

%% answer: 蜡烛应向远离透镜的方向移动，移动的距离为 $18\Ucm$
\jdanswer{
蜡烛应向远离透镜的方向移动，移动的距离为 $18\Ucm$
}

%% memo:
\memoanswer{
解：先求蜡烛的原位置。由放大率公式 $\frac{|v_{1}|}{u_{1}}=4$ 得 $v_{1}=-4u_{1}$ ①。由透镜成像公式 $\frac{1}{u_{1}}+\frac{1}{v_{1}}=\frac{1}{f}$ ②，解得 $u_{1}=\frac{3}{4}f$。再求蜡烛移动后的位置，由放大率公式得 $v_{2}=4u_{2}$ ③。由透镜成像公式 $\frac{1}{u_{2}}+\frac{1}{v_{2}}=\frac{1}{f}$ ④，解得 $u_{2}=\frac{5}{4}f$。所以蜡烛应向远离透镜的方向移动，移动的距离为 $u_{2}-u_{1}=\frac{5}{4}f-\frac{3}{4}f=\frac{1}{2}f=18\Ucm$ ⑤。

评分标准：本题 9 分。①式 2 分，②式 1 分，③式 2 分，④式 1 分，⑤式 2 分。物体移动方向正确的给 1 分。
}

\item
%% number: 23
%% paperName: 1997年全国高考第 23 题 9 分
%% typeId: 6
%% type: 计算题
%% chapter: 气体性质
%% point: 玻意耳定律
%% method: 无
%% score: 9
%% degree: 650
%% duplicateId: 0
%% body:
图中竖直圆筒是固定不动的，粗筒横截面积是细筒的 4 倍，细筒足够长。粗筒中 A、B 两轻质活塞间封有空气，气柱长 $l=20\Ucm$。活塞 A 上方的水银深 $H=10\Ucm$，两活塞与筒壁间的摩擦不计。用外力向上托住活塞 B，使之处于平衡状态，水银面与粗筒上端相平。现使活塞 B 缓慢上移，直至水银的一半被推入细筒中，求活塞 B 上移的距离。设在整个过程中气柱的温度不变，大气压强 $p_{0}$ 相当于 $75\UcmHg$。
\onepicture[num=23, showlabel=false, label=1997全国23, align=right, w=4cm]{\includesvg{figs/23}}

%% answer: $d=8\Ucm$
\jdanswer{
$d=8\Ucm$
}

%% memo:
\memoanswer{
解：在以下的计算中，都以 $\UcmHg$ 作为压强的单位。设气体初态的压强为 $p_{1}$，则有 $p_{1}=p_{0}+H$ ①。设 $S$ 为粗圆筒的横截面积，气体初态的体积 $V_{1}=Sl$。设气体末态的压强为 $p_{2}$，有 $p_{2}=p_{0}+H/2+\frac{\frac{1}{2}HS}{\frac{1}{4}S}$ ②。设末态气柱的长度为 $l'$，气体体积为 $V_{2}=Sl'$。由玻意耳定律得 $p_{1}V_{1}=p_{2}V_{2}$ ③。活塞 B 上移的距离 $d$ 为 $d=l-l'+H/2$ ④。代入数据解得 $d=8\Ucm$ ⑤。

评分标准：本题 9 分。①式 1 分，②式 2 分，③式 1 分，④式 3 分，⑤式 2 分。
}

\item
%% number: 24
%% paperName: 1997年全国高考第 24 题 11 分
%% typeId: 6
%% type: 计算题
%% chapter: 电场
%% point: 带电粒子的圆周运动
%% method: 无
%% score: 11
%% degree: 450
%% duplicateId: 0
%% body:
在方向水平的匀强电场中，一不可伸长的不导电细线的一端连着一个质量为 $m$ 的带电小球，另一端固定于 O 点。把小球拉起直至细线与场强平行，然后无初速释放。已知小球摆到最低点的另一侧，线与竖直方向的最大夹角为 $\theta$（如图）。求小球经过最低点时细线对小球的拉力。
\onepicture[num=24, showlabel=false, label=1997全国24, align=right, w=6cm]{figs/24.png}

%% answer: $T=mg\left(3-\frac{2\cos\theta}{1+\sin\theta}\right)$
\jdanswer{
$T=mg\left(3-\frac{2\cos\theta}{1+\sin\theta}\right)$
}

%% memo:
\memoanswer{
解：设细线长为 $l$，球的电量为 $q$，场强为 $E$。若电量 $q$ 为正，则场强方向在题图中向右。从释放点到左侧最高点，重力势能的减少等于电势能的增加，$mgl\cos\theta=qEl(1+\sin\theta)$ ①。若小球运动到最低点时的速度为 $v$，此时线的拉力为 $T$，由能量关系得 $\frac{1}{2}mv^{2}=mgl-qEl$ ②。由牛顿第二定律得 $T-mg=m\frac{v^{2}}{l}$ ③。由以上各式解得 $T=mg\left(3-\frac{2\cos\theta}{1+\sin\theta}\right)$ ④。

评分标准：本题 11 分。①、②式各 3 分，③式 2 分，④式 3 分。
}

\item
%% number: 25
%% paperName: 1997年全国高考第 25 题 12 分
%% typeId: 6
%% type: 计算题
%% chapter: 运动和力
%% point: 动量和能量
%% method: 无
%% score: 12
%% degree: 350
%% duplicateId: 0
%% body:
质量为 $m$ 的钢板与直立轻弹簧的上端连接，弹簧下端固定在地上。平衡时，弹簧的压缩量为 $x_{0}$，如图所示。一物块从钢板正上方距离为 $3x_{0}$ 的 A 处自由落下，打在钢板上并立刻与钢板一起向下运动，但不粘连。它们到达最低点后又向上运动。已知物块质量也为 $m$ 时，它们恰能回到 O 点。若物块质量为 $2m$，仍从 A 处自由落下，则物块与钢板回到 O 点时，还具有向上的速度。求物块向上运动到达的最高点与 O 点的距离。
\onepicture[num=25, showlabel=false, label=1997全国25, align=right, w=4.5cm]{figs/25.png}

%% answer: $l=\frac{1}{2}x_{0}$
\jdanswer{
$l=\frac{1}{2}x_{0}$
}

%% memo:
\memoanswer{
解：物块自由落下时，与钢板碰撞时的速度为 $v_{0}$，由机械能守恒得 $mg\cdot3x_{0}=\frac{1}{2}mv_{0}^{2}$，解得 $v_{0}=\sqrt{6gx_{0}}$ ①。设 $v_{1}$ 表示质量为 $m$ 的物块与钢板碰撞后一起开始向下运动的速度，因碰撞时间极短，动量守恒，$mv_{0}=2mv_{1}$ ②。刚碰完时弹簧的弹性势能为 $E_{p}$。当它们一起回到 O 点时，弹簧无形变，弹性势能为零，根据题给条件，这时物块与钢板的速度为零，由机械能守恒，$E_{p}=\frac{1}{2}\cdot2mv_{1}^{2}=2mgx_{0}$ ③。设 $v_{2}$ 表示质量为 $2m$ 的物块与钢板碰撞后开始一起向下运动的速度，则有 $2mv_{0}=3mv_{2}$ ④。刚碰完时弹簧的弹性势能为 $E_{p}'$，它们回到 O 点时，弹性势能为零，但它们仍继续向上运动，设此时速度为 $v$，则有 $E_{p}'+\frac{1}{2}\cdot3mv_{2}^{2}=3mgx_{0}+\frac{1}{2}\cdot3mv^{2}$ ⑤。在以上两种情况中，弹簧的初始压缩量都是 $x_{0}$，故有 $E_{p}'=E_{p}$ ⑥。当质量为 $2m$ 的物块与钢板一起回到 O 点时，弹簧的弹力为零，物块与钢板只受到重力作用，加速度为 $g$。一过 O 点，钢板受到弹簧向下的拉力作用，加速度大于 $g$。由于物块与钢板不粘连，物块不可能受到钢板的拉力，其加速度仍为 $g$。故在 O 点物块与钢板分离，分离后，物块以速度 $v$ 竖直上升，则由以上各式解得，物块向上运动所到最高点与 O 点的距离为 $l=\frac{v^{2}}{2g}=\frac{1}{2}x_{0}$ ⑦。故物块向上运动所到最高点与 O 点的距离为 $\frac{1}{2}x_{0}$。

评分标准：本题 12 分。①、②、③、④式各 1 分，⑤式 2 分，⑥式 3 分，得出⑦式再给 3 分。
}

\item
%% number: 26
%% paperName: 1997年全国高考第 26 题 12 分
%% typeId: 6
%% type: 计算题
%% chapter: 电场
%% point: 带电粒子的偏转
%% method: 无
%% score: 12
%% degree: 250
%% duplicateId: 0
%% body:
如图甲所示，真空室中电极 K 发出的电子（初速不计）经过 $U_{0}=1000\UV$ 的加速电场后，由小孔 S 沿两水平金属板 A、B 间的中心线射入。A、B 板长 $l=0.20\Um$，相距 $d=0.020\Um$，加在 A、B 两板间电压 $u$ 随时间 $t$ 变化的 $u-t$ 图线如图乙所示。设 A、B 间的电场可看作是均匀的，且两板外无电场。在每个电子通过电场区域的极短时间内，电场可视作恒定的。两板右侧放一记录圆筒，筒在左侧边缘与极板右端距离 $b=0.15\Um$，筒绕其竖直轴匀速转动，周期 $T=0.20\Us$，筒的周长 $s=0.20\Um$，筒能接收到通过 A、B 板的全部电子。
\begin{enumerate}
	\item
	以 $t=0$ 时（见图乙，此时 $u=0$）电子打到圆筒记录纸上的点作为 $xy$ 坐标系的原点，并取 $y$ 轴竖直向上。试计算电子打到记录纸上的最高点的 $y$ 坐标和 $x$ 坐标。（不计重力作用）
	\item
	在给出的坐标纸（图丙）上定量地画出电子打到记录纸上的点形成的图线。
\end{enumerate}
\threepicture[numA=甲, numB=乙, numC=丙, labelA=1997全国26a, labelB=1997全国26b, labelC=1997全国26c, align=right, w=5.2cm]
{figs/26a.png}
{figs/26b.png}
{figs/26c.png}

%% answer: （1）$y=2.5\Ucm$，$x_{1}=2\Ucm$，$x_{2}=12\Ucm$；（2）如图所示
\jdanswer{
\begin{enumerate}
	\item
	$y=2.5\Ucm$，$x_{1}=2\Ucm$，$x_{2}=12\Ucm$
	\item
	如图所示
\end{enumerate}
}

%% memo:
\memoanswer{
解：（1）计算电子打到记录纸上的最高点的坐标。粒子的轨迹如图所示，设 $v_{0}$ 为电子沿 A、B 板的中心线射入电场时的初速度，由动能定理得 $eU_{0}=\frac{1}{2}mv_{0}^{2}$ ①。电子在中心线方向的运动为匀速运动，设电子穿过 A、B 板的时间为 $t_{0}$，则 $l=v_{0}t_{0}$ ②。电子在垂直 A、B 板方向的运动为匀加速直线运动。对于恰能穿过 A、B 板的电子，在它通过时加在两板间的电压 $u_{c}$ 应满足 $\frac{1}{2}d=\frac{1}{2}\frac{eu_{c}}{md}t_{0}^{2}$ ③。联立①②③式解得 $u_{c}=\frac{2d^{2}}{l^{2}}U_{0}=20\UV$。此电子从 A、B 板射出时沿 $y$ 方向的分速度为 $v_{y}=\frac{eu_{c}}{md}t_{0}$ ④。此后，此电子作匀速直线运动，它打在记录纸上的点最高，设纵坐标为 $y$，根据轨迹图，由几何关系得 $\frac{y-\frac{d}{2}}{b}=\frac{v_{y}}{v_{0}}$ ⑤。由以上各式解得 $y=\frac{bd}{l}+\frac{d}{2}=2.5\Ucm$ ⑥。从题给的 $u-t$ 图线可知，加于两板电压 $u$ 的周期 $T_{0}=0.10\Us$，$u$ 的最大值 $u_{m}=100\UV$，因为 $u_{c}<u_{m}$，在一个周期 $T_{0}$ 内，只有开始的一段时间间隔 $\Delta t$ 内有电子通过 A、B 板，$\Delta t=\frac{u_{c}}{u_{m}}T_{0}$ ⑦。因为电子打在记录纸上的最高点不止一个，根据题中关于坐标原点与起始记录时刻的规定，第一个最高点的 $x$ 坐标为 $x_{1}=\frac{\Delta t}{T}s=2\Ucm$ ⑧，第二个最高点的 $x$ 坐标为 $x_{2}=\frac{\Delta t+T_{0}}{T}s=12\Ucm$ ⑨，第三个最高点的 $x$ 坐标为 $x_{3}=\frac{\Delta t+2T_{0}}{T}s=22\Ucm$。由于记录筒的周长为 $20\Ucm$，所以第三个最高点已与第一个最高点重合，即电子打到记录纸上的最高点只有两个，它们的 $x$ 坐标分别由⑧和⑨表示。坐标点表示为（$2\Ucm$，$2.5\Ucm$）、（$12\Ucm$，$2.5\Ucm$）。

\onepicture[align=right, w=6cm]{figs/26_an1.jpg}

（2）电子打到记录纸上所形成的图线，如图（2）所示。

\onepicture[align=right, w=7cm]{figs/26_an2.jpg}

评分标准：本题 12 分。第（1）问 10 分，①、②、③、④、⑤、⑥式各 1 分，⑦式 2 分，⑧、⑨式各 1 分。第（2）问 2 分，有任何错误都不给这 2 分。
}

\end{enumerate}

\end{document}
